\chapter{Visual State, Planning, and Closed-Loop Execution}
\label{ch:vla-planning}

\chapteroverview
  {Explains the optional learned-visual and task-execution stack, including its
   deterministic safety boundary and immutable deployment lifecycle.}
  {Follow the live path first---visual observation, confirmed goal, planning,
   supervised primitive, and verified effect---then use the offline sections
   for dataset, promotion, authorization, and controlled change work. Base
   simulation and direct rail control do not require this optional stack.}

\section{Layering and Safety Intent}

The optional research stack deliberately separates perception, language,
planning, supervision, and actuation. No learned component publishes a typed
rail command. The deterministic supervisor is the only component in this stack
that does so.

\begin{longtable}{L{0.16\textwidth} L{0.29\textwidth} L{0.19\textwidth} L{0.27\textwidth}}
\caption{Closed-loop runtime stages and authority boundaries}
\label{tab:closed-loop-stack}\\
\toprule
\tableheader Stage & Exchange & Authority/gate & Source file(s) \\
\midrule
\endfirsthead
\toprule
\tableheader Stage & Exchange & Authority/gate & Source file(s) \\
\midrule
\endhead
Observation readiness & Gazebo effect $\rightarrow$ paired images/presence $\rightarrow$ accepted V4 state & Promoted model, validator, presence/fusion readiness & \file{mfja_robot_control_config/scripts/room_315_visual_state_inference_node.py}; \file{mfja_robot_control_config/scripts/room_315_visual_runtime_validation.py} \\ \rowseparator
Goal intake & Constrained English $\rightarrow$ confirmed \code{TaskGoal} & Parser, validation, explicit operator confirmation & \file{mfja_robot_control_config/scripts/room_315_task_goal_cli.py}; \file{mfja_robot_control_config/scripts/room_315_task_goal_builder.py} \\ \rowseparator
Grounding & Fresh accepted observation $\rightarrow$ concrete task facts & Task gateway schema/checkpoint/freshness gates & \file{mfja_robot_control_config/scripts/room_315_task_execution.py} \\ \rowseparator
Planning & Fresh PDDL problem $\leftrightarrow$ PlanSys2 plan & Runtime domain plus whole-plan validation & \file{mfja_robot_control_config/config/room_315_vla/pddl/domain_room315_runtime.pddl}; \file{mfja_robot_control_config/scripts/room_315_pddl_plan_translator.py} \\ \rowseparator
Primitive execution & First translated primitive $\leftrightarrow$ supervisor status & Gateway authorization and deterministic supervisor & \file{mfja_robot_control_config/scripts/room_315_task_execution.py}; \file{mfja_robot_control_config/scripts/room_315_vla_supervisor.py} \\ \rowseparator
Typed actuation & Shuttle/switch/stopper command $\leftrightarrow$ actual state & Rail controller and typed interfaces & \file{mfja_robot_control_config/scripts/room_315_kinematic_shuttle_node.py}; \file{mfja_rail_interfaces/CMakeLists.txt} \\ \rowseparator
Effect verification & Newer visual/controller evidence $\rightarrow$ success or replan & Gateway effect predicates and bounded retries; the loop returns to observation readiness & \file{mfja_robot_control_config/scripts/room_315_task_execution.py}; \file{mfja_robot_control_config/scripts/room_315_visual_state_inference_node.py} \\
\bottomrule
\end{longtable}

\section{V4 Visual-State Runtime}

Source files: \file{mfja_robot_control_config/launch/room_315_visual_state_runtime.launch.py},
\file{mfja_robot_control_config/scripts/room_315_visual_state_inference_node.py},
\file{mfja_robot_control_config/scripts/room_315_visual_runtime_v4.py}, and
\file{mfja_robot_control_config/config/room_315_vla/visual_state_runtime.yaml}.

\file{room_315_visual_state_runtime.launch.py} starts the separate learned
runtime. Only generation V4 is admitted to the active path. A promoted runtime
must supply the exact manifest and checkpoint identities expected by config;
missing files, wrong hashes, wrong schema, stale inputs, rejected predictions,
or unsafe supervisor status fail closed.

The inference node consumes:

\begin{itemize}
  \item synchronized left/right RGB images;
  \item each rail's \code{ShuttleState} only for identity presence, timestamp,
        receive time, and known-empty heartbeat; and
  \item deterministic supervisor safety status.
\end{itemize}

It must not use controller segment, arc length, world pose, slot, payload state,
or privileged identity-tracker geometry as learned inputs. This separation
prevents evaluation leakage. Presence requires both rail heartbeats, a 0.5 s
warm-up, and freshness no older than 1.0 s under current defaults. A blank-name
heartbeat is known-empty; no heartbeat is unknown/stale.

The node publishes typed raw prediction, validation, and accepted observation
streams plus diagnostic JSON/model output. The default runtime mode is
\code{active}, but \code{dry_run_state_fusion=true} and
\code{plansys2_update_enabled=false}; therefore merely starting visual runtime
does not mutate PlanSys2 knowledge. The dry-run flag gates external predicate
writes; it does not suppress internal fusion or accepted-observation
publication. \code{shadow} mode writes under
\rosname{/room_315/visual_state/shadow_v4/...} and is observation-only.

\section{Visual Observation Contract}

Source files: \file{mfja_rail_interfaces/msg/VisualStateObservation.msg},
\file{mfja_rail_interfaces/msg/VisualShuttleState.msg},
\file{mfja_robot_control_config/scripts/room_315_visual_runtime_validation.py},
and \file{mfja_robot_control_config/scripts/room_315_task_execution.py}.

For every canonical identity \code{L1--L4,R1--R4}, the model/validator may
report presence, side/block, bounding box, normalized/metric segment position,
loaded state, and confidence. The enclosing observation records source image
stamps, schema/checkpoint, acceptance/stabilization/readiness flags, rejection
reasons, clamps, latency, and counters.

Consumers must apply all of the following rather than checking only
\code{accepted}:

\begin{enumerate}
  \item message and both source images are fresh;
  \item schema and checkpoint hash exactly match the authorized values;
  \item model/input/presence/fusion readiness is true;
  \item the observation is accepted and not stale;
  \item facts required by the requested task are valid and sufficiently
        confident; and
  \item a newer observation follows any commanded action before effect success.
\end{enumerate}

\code{stabilized} reports that the optional temporal filter was applied. It is
normally false when \code{temporal_filter_enabled=false}, including on valid
accepted observations, and the task gateway does not require it. If a promoted
runtime enables temporal filtering, validate that filter's authorized contract
rather than imposing \code{stabilized=true} on every deployment.

\section{English TaskGoal Understanding}

Source files: \file{mfja_robot_control_config/scripts/room_315_task_goal_cli.py},
\file{mfja_robot_control_config/scripts/room_315_task_goal_dialogue.py},
\file{mfja_robot_control_config/scripts/room_315_task_goal_parsers.py}, and
\file{mfja_robot_control_config/config/room_315_vla/task_goal_understanding.yaml}.

The language layer combines a deterministic parser with an optional local
semantic model. The deterministic path remains available without the large
model. It produces a draft, asks targeted clarification for ambiguity, and
requires an explicit yes/confirm step before publishing JSON to
\rosname{/room_315/task_goal}.

Supported intent is deliberately narrow: one transport or inspection request,
explicit/any/nearest identity where valid, loaded/empty/any payload
qualification, and an exact slot or supported station target. Compound plans,
swaps, ``move all'', load/unload operations, arbitrary coordinates, deadlines,
and speed-optimization objectives are rejected rather than guessed.

The optional model requires the separately managed
\code{llama-cpp-python==0.3.16} environment and a roughly 1.04 GiB GGUF file.
\file{mfja_robot_control_config/scripts/setup_room315_intent_model.py} can
download/install dependencies and may
use user-site or \code{--break-system-packages}; prefer the documented isolated
virtual environment and inspect its planned actions before running it.

\section{PDDL Authorities and Translation}

Source files: \file{mfja_robot_control_config/config/room_315_vla/pddl/domain_room315_runtime.pddl},
\file{mfja_robot_control_config/config/room_315_vla/pddl/domain_room315.pddl},
\file{mfja_robot_control_config/scripts/room_315_pddl_plan_translator.py}, and
\file{mfja_robot_control_config/scripts/room_315_pddl_validation_gate.py}.

\file{domain_room315_runtime.pddl} is the live execution domain.
\file{domain_room315.pddl} is the broader expert/data-generation domain. Do not
replace the runtime file with the research domain merely because it contains
more actions.

The current runtime domain defines these 21 actions:

\begin{lstlisting}
prepare_switches                 open_stoppers
restore_normal_route             move_shuttle_to_slot
prepare_topology_route           move_shuttle_from_segment_to_slot
prepare_slot_topology_route      move_shuttle_via_topology_to_slot
begin_route_clearance            relocate_blocker_to_interior
stage_selected_to_interior       finish_route_clearance
begin_segment_route_clearance    relocate_segment_blocker_to_interior
stage_selected_segment_to_interior
finish_segment_route_clearance   pause_route_clearance
stop_shuttle                     finish_task
finish_candidate_task            inspect_state
\end{lstlisting}

The expert and runtime domains currently expose the same action names, while
their predicates/preconditions remain separate versioned contracts. The
checked runtime file is always the machine authority. The translator reduces
the supported plan surface to
deterministic primitives: \code{SET_SWITCHES}, \code{SET_STOPPERS},
\code{SHUTTLE_ON}, \code{STOP_NOW}, and \code{DONE}. The entire returned plan
is parsed and validated before its first primitive is sent.

\section{Closed-Loop Task Gateway}

Source files: \file{mfja_robot_control_config/scripts/room_315_task_execution_node.py},
\file{mfja_robot_control_config/scripts/room_315_task_execution.py},
\file{mfja_robot_control_config/scripts/room_315_task_execution_config.py}, and
\file{mfja_robot_control_config/config/room_315_vla/task_execution_runtime.yaml}.

\file{room_315_task_execution_node.py} owns the host singleton lock and combines
the accepted visual observation, supervisor status, and both typed sensor
feeds. It ignores privileged \code{segment}, \code{s}, and \code{s_ratio}
fields in sensor readings; only active state, sensor name, and shuttle identity
may anchor exact DZI occupancy.

For each confirmed goal the gateway:

\begin{enumerate}
  \item verifies execution enablement, Gazebo-only policy, immutable task
        authorization, promotion manifest, schema, and checkpoint hash;
  \item waits for fresh visual/supervisor/sensor inputs;
  \item grounds identity and payload constraints; payload requests require five
        agreeing accepted frames within at most 15 observations;
  \item generates a new problem and calls \rosname{/planner/get_plan};
  \item validates/translates the complete response;
  \item executes only the first primitive through the supervisor;
  \item waits for a newer visual/controller effect, then replans; and
  \item succeeds only when the exact goal holds, or publishes \code{stop_all}
        and aborts on bounded failure.
\end{enumerate}

An exact final slot requires both identity-bearing DZI occupancy and explicit
controller \state{DISABLED}. Current defaults include 20 s planner timeout,
1.5 s visual/supervisor freshness, 1.0 s sensor freshness, 2.0 s observation
wait, two confirming sensor frames, 0.2 m/s motion, 32 maximum steps, eight
replans, three unknown-effect retries, 5 s supervisor timeout, 30 s ordinary
effect timeout, and 60 s clearance timeout. Long-route arrival uses
$1.25\times$ estimated travel time plus 5 s margin.

\section{Primitive Supervisor}

Source files: \file{mfja_robot_control_config/scripts/room_315_vla_supervisor.py},
\file{mfja_robot_control_config/launch/room_315_vla_supervisor.launch.py}, and
\file{mfja_robot_control_config/config/room_315_vla/vla_supervisor.yaml}.

\file{room_315_vla_supervisor.py} is deterministic. It consumes JSON primitives
on \rosname{/room_315/vla/command}, reads typed rail/device state and camera
availability, and publishes typed shuttle, switch, and stopper commands. It
reports structured status on \rosname{/room_315/vla/status}.

Important invariants are:

\begin{itemize}
  \item arbitrary learned action vectors are rejected;
  \item identities, sides, slots, switch/stopper states, speeds, readiness, and
        incompatible simultaneous actions are validated before publication;
  \item \code{stop_all} is the normal fail-safe primitive;
  \item an emergency stop closes stoppers and prevents continued actuation;
        clearing the Boolean emergency-stop input does not automatically
        restart previously moving shuttles; and
  \item bounded recovery is attempted at most twice by default, after which the
        supervisor enters emergency-stop/manual-recovery behavior.
\end{itemize}

The supervisor's legacy position-feedback subscriptions are not current state
authority. Typed shuttle, switch, stopper, and \code{sensors/feedback} topics
are the live feedback surface.

\section{Dataset Lifecycle and Leakage Control}

Source files: \file{mfja_robot_control_config/scripts/room_315_vla_dataset_recorder.py},
\file{mfja_robot_control_config/scripts/room_315_visual_state_dataset.py},
\file{mfja_robot_control_config/scripts/room_315_grouped_visual_splits.py}, and
\file{mfja_robot_control_config/scripts/room_315_visual_dataset_audit.py}.

A recorded episode contains \file{images/<camera>/}, \file{data.jsonl},
\file{events.jsonl}, \file{episode.json}, and, after validation,
\file{validation.json}. The supported lifecycle is summarized below.

\begin{longtable}{L{0.20\textwidth} L{0.36\textwidth} L{0.35\textwidth}}
\caption{Dataset-to-candidate lifecycle and owning tools}
\label{tab:visual-dataset-lifecycle}\\
\toprule
\tableheader Stage & Required transition/output & Source file(s) \\
\midrule
\endfirsthead
\toprule
\tableheader Stage & Required transition/output & Source file(s) \\
\midrule
\endhead
Capture & Scenario manifest $\rightarrow$ guarded episode images/events/data & \file{mfja_robot_control_config/scripts/room_315_vla_dataset_recorder.py}; \file{mfja_robot_control_config/scripts/room_315_visual_state_capture.py} \\ \rowseparator
Validate/label & Episode $\rightarrow$ validation record and permitted labels/export & \file{mfja_robot_control_config/scripts/room_315_visual_state_dataset.py}; \file{mfja_robot_control_config/scripts/room_315_visual_label_exporter.py} \\ \rowseparator
Split/audit & Related captures remain grouped; leakage and role constraints are audited & \file{mfja_robot_control_config/scripts/room_315_grouped_visual_splits.py}; \file{mfja_robot_control_config/scripts/room_315_visual_dataset_audit.py} \\ \rowseparator
Train/evaluate & Frozen training selection $\rightarrow$ sealed evaluation evidence & \file{mfja_robot_control_config/scripts/room_315_visual_training_v4.py}; \file{mfja_robot_control_config/scripts/room_315_visual_v4_final_test.py} \\ \rowseparator
Package & Verified model/config/evidence $\rightarrow$ immutable candidate & \file{mfja_robot_control_config/scripts/room_315_build_runtime_candidate_v4.py} \\
\bottomrule
\end{longtable}

Privileged identity tracks and controller coordinates may be used as label or
evaluation oracle data only when the applicable contract permits them. They
must remain separate from model inputs. Split tools group related captures so
variants of the same scene do not cross train/validation/test boundaries.

The frozen V4 training contract uses a 50/50 legacy/hard-data mixture, a
split-rail ResNet-18 model at $320\times240$, 12 epochs, and seed
31520260811. The associated evidence records these frozen roles.

The tracked role/count authorities are
\file{mfja_robot_control_config/config/room_315_vla/visual_state_training_v4.json},
\file{mfja_robot_control_config/config/room_315_vla/visual_state_final_test_v4.json},
and \file{mfja_robot_control_config/config/room_315_vla/visual_state_final_test_v4_coverage_extension.json}.

\begin{longtable}{L{0.32\textwidth} L{0.18\textwidth} L{0.40\textwidth}}
\caption{Frozen V4 dataset roles, counts, and permitted use}
\label{tab:v4-dataset-roles}\\
\toprule
\tableheader Dataset role & Count & Use \\
\midrule
\endfirsthead
\toprule
\tableheader Dataset role & Count & Use \\
\midrule
\endhead
Legacy training & 1,528 & Historical training component \\ \rowseparator
Historical predecessor validation & 256 & Provenance/comparison; not reusable as a new final test \\ \rowseparator
Consumed historical test & 256 & Already consumed; never tune against it again \\ \rowseparator
Hard training & 4,000 & V4 hard-case training component \\ \rowseparator
V4 validation & 512 & Model/epoch selection under frozen rules \\ \rowseparator
Canary & 256 & Guarded pre-deployment check \\ \rowseparator
Sealed final-test coverage & 1,040 & Frozen extension of the 1,024-scene base final-test contract \\
\bottomrule
\end{longtable}

Treat those values as the current contract, not a target to reproduce by
silently adding files. Any changed dataset requires a new version, immutable
manifest, provenance, split audit, and qualification decision.

\section{Promotion and Authorization}

Source files: \file{mfja_robot_control_config/scripts/room_315_build_runtime_candidate_v4.py},
\file{mfja_robot_control_config/scripts/room_315_promote_runtime_v4.py},
\file{mfja_robot_control_config/scripts/room_315_authorize_runtime_v4_closed_loop.py},
\file{mfja_robot_control_config/scripts/room_315_visual_runtime_v4.py}, and
\file{mfja_robot_control_config/scripts/room_315_task_execution_config.py}.

The intended chain is:

\begin{enumerate}
  \item build an immutable candidate with dataset, config, code, and checkpoint
        hashes;
  \item run validation and CPU/runtime compatibility checks;
  \item deploy in shadow mode and record acceptance/fault evidence;
  \item manually promote the exact candidate, producing a signed/hashed
        promotion manifest;
  \item separately authorize closed-loop task execution for that exact
        manifest/checkpoint/code contract; and
  \item point runtime config at those immutable files without editing the
        candidate in place.
\end{enumerate}

Promotion permits inference; task authorization permits a specific Gazebo
closed-loop configuration. They are intentionally separate. Rollback selects a
previous complete authorized bundle and restores all matching hashes/config;
never copy only an older checkpoint into a newer manifest.

Task authorization is verified only when \code{execution_enabled=true}; a
disabled gateway requires no authorization artifact. For enabled execution,
the authorization JSON, promotion manifest, manual decision, and every bound
artifact must be regular files whose final path component is not a symbolic
link and whose Unix mode has no write bits. Hash agreement does not replace
these immutability checks.

\begin{warningbox}[Clean-clone limitation]
The active checked-in V4 visual and task-execution YAML files reference
absolute \file{/home/tiago/...} promotion and authorization artifacts that are
not stored in Git. A clean clone can run simulation, rails, cameras, recording,
and deterministic supervision, but active learned inference and authorized
execution require the complete external bundle to be restored or a new bundle
to be qualified. Do not weaken hash checks to make startup succeed.
\end{warningbox}

\section{Changing the Learned/Planning Stack}

The coupled change surface is implemented and tested across
\file{mfja_robot_control_config/scripts/room_315_visual_model_v4.py},
\file{mfja_robot_control_config/scripts/room_315_visual_runtime_validation.py},
\file{mfja_robot_control_config/scripts/room_315_pddl_plan_translator.py},
\file{mfja_robot_control_config/scripts/room_315_task_execution.py}, and
\file{mfja_robot_control_config/test/test_room315_task_execution.py}.

\begin{procedurebox}[Safe change sequence]
\begin{enumerate}
  \item Version the contract; never edit a sealed dataset, checkpoint, manifest,
        authorization, or final-test result in place.
  \item Keep controller/oracle data out of learned inputs and record the exact
        feature boundary in tests.
  \item Update model, validator, typed observation schema, state builder, PDDL
        generator/domain, translator, executive, config, and fault tests as one
        reviewed dependency set.
  \item Train/select only on allowed splits; preserve the final test until the
        decision rule is frozen.
  \item Build a new immutable candidate, run shadow/acceptance/fault campaigns,
        promote manually, and issue a separate execution authorization.
  \item Retain a complete prior bundle and rehearse rollback before enabling
        execution.
\end{enumerate}
\end{procedurebox}
